\makeatletter
    \setlength{\absparsep}{18pt}
\makeatother

\begin{center}
    \textbf{RESUMO}
\end{center}
\vspace{.5cm}

\noindent
As cargas dinâmicas induzidas por ondas em Plataformas de Pernas Tracionadas representam uma ameaça à integridade estrutural, sobretudo pela transmissão de vibrações ao longo dos tendões. Neste contexto, este estudo investiga o uso de configurações periódicas nesses elementos como estratégia passiva para mitigação de vibrações. A partir da teoria de vigas de Euler-Bernoulli, desenvolveu-se uma rotina numérica baseada no Método dos Elementos Finitos, implementada em MATLAB e validada com a teoria analítica, para analisar a formação de faixas de atenuação de frequência, onde a propagação de ondas é suprimida. Duas configurações periódicas são avaliadas: vigas escalonadas (periodicidade geométrica) e vigas bi-materiais (periodicidade de material). Uma análise da formação de zonas de atenuação é realizada por meio de um parâmetro de saturação periódica. Os resultados mostram que, embora ambas as configurações sejam capazes de gerar band gaps, as vigas escalonadas apresentam maior capacidade de sintonização e produzem bandas de atenuação mais amplas. Assim, a otimização da periodicidade geométrica destaca-se como uma abordagem de projeto robusta para a filtragem de frequências de vibração críticas em sistemas marítimos.

\vspace{18pt}\noindent\textbf{Palavras-chave}: Plataformas de Pernas Tracionadas, Vigas Periódicas, Mitigação de Vibrações, Faixas de Atenuação, Método dos Elementos Finitos.